%{{{preamble
\documentclass[a4paper,10pt]{article} \usepackage{anysize}
\marginsize{2cm}{2cm}{1cm}{1cm}
%\textwidth 6.0in \textheight = 664pt
\usepackage{xltxtra}
\usepackage{graphicx}
\usepackage{color} 
\usepackage{xgreek} 
\usepackage{fancyvrb}
\usepackage{minted}
\usepackage{listings}
\usepackage{enumitem} 
\usepackage{framed} 
\usepackage{relsize}
\usepackage{float} 
\setmainfont[Mapping=TeX-text]{Kerkis}
\renewcommand{\labelenumi}{\alph{enumi}}
\renewcommand{\labelenumii}{(\roman{enumii})}

%}}}
\begin{document}

\begin{titlepage}
\begin{center}
\begin{figure}[t] 
     \includegraphics[scale=0.7]{title/ntua_logo}
\end{figure}
\begin{LARGE}\textbf{ΕΘΝΙΚΟ ΜΕΤΣΟΒΙΟ ΠΟΛΥΤΕΧΝΕΙΟ\\}\end{LARGE}
\vspace{5cm}
\begin{Large}
ΣΧΟΛΗ ΗΜ\&ΜΥ\\
Αλγόριθμοι και Πολυπλοκότητα\\
1\textsuperscript{η} Σειρά Γραπτών Ασκήσεων\\
Ακ. έτος 2010-2011\\
\end{Large}
\vspace{10cm}
\begin{tabular}{l r}
\Large{Λύρας Γρηγόρης}&
\large{Α.Μ.: 03109687}\\
\end{tabular}\\
\vspace{1cm}

\vfill
\large\today\\
\end{center}
\end{titlepage}


\section{Ασυμπτωτικός συμβολισμός, Αναδρομικές
Σχέσεις} \setcounter{section}{1}
\begin{enumerate}
%{{{ Ταξινόμηση
\item \textit{Να ταξινομήσετε τις παρακάτω συναρτήσεις σε αύξουσα σειρά τάξης
μεγέθους, να βρείτε δηλαδή μια διάταξη $g_1,g_2,g_3,...$ τέτοια ώστε $g_1 =
O(g_2),g_2=O(g_3),$ κοκ. Σε αυτή τη διάταξη, να επισημάνετε τις συναρτησεις
που έχουν ίδια τάξη μεγέθους.}
\begin{enumerate}
\item $\log {n^3}$
\item $\sqrt{n}*\log^{50}{n} $
\item $\frac{n}{\log{\log{n}} } $
\item $\log{n!} $
\item $n*\log^{10}{n} $
\item $n^{1.01} $
\item $5^{\log_{2}{n}} $
\item $\sum_{k=1}^n{k^5} $
\item $\log^{\log{n}}{n} = n^{\log{\log{n}} } $
\item $2^{\log_2^4{n}} $
\item $\log^{\sqrt{n}}{n} $
\item $e^{\frac{n}{\ln{n}} } $
\item $n*3^n $
\item $2^{2*n} $
\item $\sqrt{n!} $
\end{enumerate}
%}}}
%{{{ Τάξη Μεγέθους
\item \textit{Να υπολογίσετε την τάξη μεγέθους $\Theta$ των λύσεων των παρακάτω αναδρομικών
σχέσεων. Για όλες τις σχέσεις, να θεωρήσετε ότι $T(1)=\Theta(1)$}
\begin{enumerate}
\item $T(n)=5*T(n/7)+n*\log(n)$
$\Rightarrow n^{\log _7 {5}} = n^{0.827}$
$\Rightarrow n^{0.827} < n*\log{n}$\\
$\Rightarrow T(n) \in \Theta(n*\log{n})$

\item $T(n) = 4*T(n/5)+n/\log^2{n}$
$\Rightarrow n^{\log_5 {4}} = n^{0.861}$
$\Rightarrow n^{0.861} > n/\log^2{n}$\\
$\Rightarrow T(n) \in \Theta(n/\log^2{n})$

\item $T(n) = T(n/3) +3*T(n/7)+n$\\
$\Rightarrow T(n) \in O(n)$

\item $T(n) = 6*T(n/6)+n$
$\Rightarrow n^{\log _6 {6}} = n$\\
$\Rightarrow T(n) \in \Theta(n*\log {n})$

\item $T(n) = T(n/3) + T(2n/3) + n$\\
$\Rightarrow T(n) \in \Theta(n*\log{n})$

\item $T(n) = 16*T(n/4) + n^3*\log^2{n}$
$\Rightarrow n^{\log _4{16}} = n^2$
$\Rightarrow n^2 < n^3*\log^2{n}$\\
$\Rightarrow T(n) \in \Theta(n^3*\log^2{n})$

\item $T(n) = T(\sqrt(n)) + \Theta(\log{\log{n}})$
$(k=\log{n},f(k)=T(n))$\\
$\Rightarrow
T(\sqrt{n})=f(\log{\sqrt{n}})=f(\frac{1}{2}*\log{n})=f(\frac{k}{2}) $\\
$\Rightarrow f(k)=f(\frac{k}{2})+\Theta(\log{k})$\\
$\Rightarrow \Theta(\frac{1}{2}*\log^2{k}*(\log{k}+1))$

\item $T(n) = T(n-3) + \log{n}$\\
$\Rightarrow T(n) \in \Theta(n*logn)$

\end{enumerate}
%}}}
\end{enumerate}
%{{{Ταξινόμηση σε Πίνακα με Πολλά Ίδια Στοιχεία
\pagebreak
\section{Ταξινόμηση σε Πίνακα με Πολλά Ίδια Στοιχεία}
\textit{Έστω πίνακας ακεραίων $A[1...n]$ που χαρακτηρίζεται από πολλές
πολλαλπές εμφανίσεις των στοιχείων του. Συγκεκριμένα, θεωρούμε ότι το πλήθος
των διαφορετικών στοιχείων του Α είναι μόλις πολυλογαριθμικό (δηλ.
$O(\log^d{n})$, για κάποια σταθερά $d \geq 1$). Να διατυπώσετε έναν συγκριτικό
αλγόριθμο που ταξινομεί τον πίνακα Α σε χρόνο $O(n*\log{\log{n}})$. Γιατί δεν
ισχύει το κάτω φράγμα του $\Omega({n*\log{n}})$ σε αυτή την περίπτωση;} \vspace{0.4cm}

Για να κάνουμε ταξινόμηση σε ένα τέτοιο πίνακα σε πρώτο στάδιο μετράμε τα
στοιχεία που είναι ίδια σε κάθε πέρασμα και τα βγάζουμε από τον πίνακα. Για να
ολοκληρώσουμε αυτή τη διαδικασία και με δεδομένο πως το πλήθος των
διαφορετικών στοιχείων είναι μόλις $\log^d{n}$ προκύπτει ο αναδρομικός τύπος
$T(n)=T(n*\frac{m-1}{m})+f(n),f(n) \in O(n),m=\log^d{n}$. Επιλύοντας με χρήση
του \textit{Master Theorem} καταλήγουμε σε πολυπλοκότητα $O(n)$ για το πρώτο
βήμα. Έπειτα με mergesort ταξινομούμε τα στοιχεία που έχουμε και τα
επεκτείνουμε ανάλογα με το πλήθος του καθενός έτσι ώστε το τελικό αποτέλεσμα 
να είναι και πάλι μήκους $n$. Η διαδικασία ολοκληρώνεται σε άλλο
$\log^d{n}*\log{\log^d{n}}=d*\log^d{n}*\log{\log{n}}$. Προφανώς το κύριο
κομμάτι του χρόνου  καταναλώνεται στο δεύτερο στάδιο. Συνεπώς η πολυπλοκότητα
του αλγορίθμου είναι $\Theta{\log^d{n}*\log{\log{n}} }$ άρα και
$O(n*\log{\log{n}})$.
%}}}
%{{{Δυαδική Αναζήτηση
\section{Δυαδική Αναζήτηση}
\begin{enumerate}
\item \textit{Μας δίνεται πίνακας Α με n ακέραιους ταξινομιμένους σε αύξουσα
σειρά, αλλά \emph{δεν γνωρίζουμε την τιμή του n}. Ο τύπος δεδομένων του Α έχει
υλοποιηθεί ώστε να επιστρέφει το μήνυμα σφάλματος $\inf$ κάθε φορά που
επιχειρείται προσπέλαση στοιχείου $A[i]$ με $i > n$. Να διατυπώσετε αλγόριθμο
με χρόνο εκτέλεσης $O(\log{n})$, που δέχεται ως είσοδο έναν ακέραιο $x$, και
βρίσκει μια θέση του πίνακα Α που περιέχει τον $x$, αν υπάρχει τέτοια θέση.}

Ξεκινάμε με ένα στοιχείο. Και συγκρίνουμε με τον $x$. Όσο ο $x$ είναι
μεγαλύτερος από αυτόν διπλασιάζουμε τον αριθμό των στοιχείων και ελέγχουμε
πάλι με τον τελευταίο. Μόλις φτάσουμε σε μεγαλύτερο αριθμό από τον $x$ έστω
στη θέση $k$ και εφαρμόζουμε κλασσική δυαδική αναζήτηση στο τμήμα $A[1...k]$. 
Μέγιστος αριθμός επαναλήψεων είναι $O(\log{k})$.
\item \textit{Έστω δύο ταξινομημένοι πίνακες ακεραίων $A[1...n]$ και
$B[1...n]$. Να διατυπώσετε αλγόριθμο με χρόνο εκτέλεσης $O(\log{k})$ που
υπολογίζει το k-οστό μικρότερο στοιχείο της ένωσης των πινάκων Α και Β. Για
ευκολία, μπορείτε να υποθέσετε ότι όλα τα στοιχεία των Α και Β είναι
διαφορετικά.}

Παίρνω τα k πρώτα στοιχεία από κάθε πίνακα. Τα χωρίζω στη μέση και
συγκρίνω τα $A[k/2]$ με $B[k/2+1]$ και $A[k/2+1]$ με $B[k/2]$. Αν
$A[k/2]<B[k/2+1]$ και $B[k/2]<A[k/2+1]$ τότε το k-οστό στοιχείο είναι το
$max\{A[k/2],B[k/2]\}$. Αν $A[k/2]<B[k/2+1]$ και $B[k/2]>A[k/2+1]$ τότε
επαναλαμβάνω χρησιμοποιώντας τα στοιχεία $[A[k/2+1]..A[k]]$ και
$[B[1]..B[k/2]]$.
\end{enumerate}
%}}}
%{{{Συλλογή Comics
\pagebreak
\section{Συλλογή Comics}
\textit{Ένας φανατικός συλλέκτης comics θυμάται ότι από τη συλλογή του
αγαπημένου του geek comic έχει χάσει ένα τεύχος, αλλά δεν θυμάται ποιο.
Χρειάζεται να εντιπίσει το τεύχος αυτό, γιατί όλα τα τεύχη της σειράς
πωλουνται στο Internet σε τιμή ευκαιρίας, σε μια δημοπρασία που λήγει άμεσα.}

\textit{Ένας από τους λόγους που η συγκεκριμένη σειρά comics είναι η αγαπημένη του
έχει να κάνει με τον ιδιόρρυθμο τρόπο αρίθμησης των τευχών: η αρίθμηση των
τευχών είναι δυαδική, ξεκινά από το $0...0$ και ολοκληρώνεται στο $1...1$.
Κάθε τεύχος έχει k σελίδες, το πλήθος των τευχών είναι $n=2^k$, και τα bits
του αριθμού κάθε τεύχους είναι γραμμενα από ένα σε κάθε σελίδα. Έτσι το
μοναδικό που μπορεί να κάνει ο συλλέκτης για να εντοπίσει το τεύχος που λείπει
είναι να εξετάζει κάθε φορά το i-οστό bit της αρίθμησης ενός τεύχους j.}

\textit{Ο συλλέκτης χρειάζεται λοιπόν έναν αλγόριθμο που εντοπίζει το τεύχος που
λείπει με όσο το δυνατόν λιγότερες ερωτήσεις της μορφής "ποιο είναι το i-οστό
bit αρίθμηση του τεύχους j;". Μπορείτε να βοηθήσετε τον συλλέκτη διατυπώνοντας
έναν αλγοριθμο που χρειάζεται το πολύ $2*n$ τέτοιες ερωτήσεις;}\vspace{0.4cm}

Ζητάμε πρώτα το MSB για όλα τα τεύχη. Έτσι τα χωρίζουμε σε $\frac{n}{2}$ και
$\frac{n}{2}-1$. Προφανώς αυτό που λείπει είναι στα λιγότερα. Εφαρμόζουμε πάλι
κάνοντας $\frac{n}{2}$ ερωτήσεις και κρατάμε πάλι το υποσύνολο με το μικρότερο
πλήθος. Αναδρομικά θα καταλήξουμε στο τεύχος που λείπει από τη συλλογή έχοντας
ρωτήσει $n+\frac{n}{2}+\frac{n}{4}+..+1=2*n*(1-\frac{1}{2^n}) < 2*n$ φορές.
%}}}
%{{{ Πολυκατοικίες χωρίς Θέα
\section{Πολυκατοικίες χωρίς Θέα}
\vspace{0.4cm}
Ξεκινάμε με τον από τον $A[1]$ θέτοντας το $B[1]=0$ εφόσον δεν υπάρχει κάποιος
δυτικότερα από αυτόν. Προχωράμε στον επόμενο $A[2]$ και συγκρίνουμε το ύψος
του με τον $A[1]$. Αν $A[2]<A[1]$ τότε θέτουμε $B[2]=1$ αλλιώς συγκρίνουμε το
ύψος $A[2]$ με το ύψος του στόχου του $A[1]$. Αν πάλι δεν βρούμε ψηλότερο
τρέχουμε πάλι για το στόχο του στόχου αναδρομικά μέχρι να φτάσουμε σε 0
ή σε κάποιον ψηλότερο του $A[2]$ και τον θέτουμε ως στόχο στο πεδίο $B[2]$.
Επαναλαμβάνουμε μέχρι $n$. Μέγιστος αριθμός συγκρίσεων είναι $2n$.
%}}}

\end{document}
